%% kapa: exam=1 round=33 year=2022 period=2 form=A subject=accounting subject_ko=회계학 seq=14 number=54
%% source: Q-Net 감정평가사 기출문제 (cst00309) artlSeq=5209997 file=raw/1008/5209997/2022년도 제33회 감정평가사 1차시험 2교시(감정평가 관계법규, 회계학).pdf
%% conversion: kordoc 4.24.51 -> tools/gen1.py (auto); review=proofread-sonnet; flags=none
\begin{question}{54}
충당부채, 우발부채 및 우발자산에 관한 설명으로 옳지 않은 것은?
\begin{choices}
  \item 충당부채는 부채로 인식하는 반면, 우발부채는 부채로 인식하지 아니한다.
  \item 충당부채로 인식하는 금액은 현재의무를 보고기간 말에 이행하기 위하여 필요한 지출에 대한 최선의 추정치이어야 한다.
  \item 충당부채에 대한 최선의 추정치를 구할 때에는 관련된 여러 사건과 상황에 따르는 불가피한 위험과 불확실성을 고려한다.
  \item 예상되는 자산 처분이익은 충당부채를 생기게 한 사건과 밀접하게 관련되어 있다고 하더라도 충당부채를 측정함에 있어 고려하지 아니한다.
  \item 충당부채는 충당부채의 법인세효과와 그 변동을 고려하여 세후 금액으로 측정한다.
\end{choices}
\end{question}
\answer{5}
